\section{dp-greedy 深度题卡}

\subsection{动态规划与状态定义}
下面只在这里讲一次共性：先写暴力基线，定位最贵的重复工作，再选择能维护不变量的数据结构。每道题卡只分析本题自己的建模、瓶颈、边界和变体，复习时不要把家族共性当成答案。

\begin{principlebox}{图示：动态规划与状态定义推导链}
\begin{tabular}{@{}p{0.18\linewidth}p{0.74\linewidth}@{}}
暴力基线 & 枚举候选、完整模拟或逐个检查，先保证覆盖全部可能答案。\\
瓶颈定位 & 慢点不是递归写法，而是相同状态被反复计算。若一个状态的未来只取决于少量参数，就应保存结果并按依赖顺序计算。\\
结构选择 & 状态定义要包含足够信息但不能过度膨胀。线性 DP 看最后一步选择，二维 DP 看两个前缀或网格位置，背包看物品阶段和容量，区间 DP 看最后合并点。\\
不变量 & 填表不变量是：计算当前状态时，所有依赖状态已经算完且语义一致。滚动数组必须解释当前格读到的是上一轮还是本轮。\\
代码落地 & 数组维度和初始化要服务于状态含义。不可达用大正数或负数哨兵时要避免溢出。方案数可能超过 int 时改用 \code{long long} 或按题目取模。\\
\end{tabular}
\end{principlebox}

\subsection{LeetCode 53 最大子数组和}

\textbf{题目模型。} 状态是必须以当前位置结尾的最大连续和。

\textbf{推导重点。} 若前一个后缀为负，接上它只会拖累当前元素，因此可以重新开始。

\textbf{边界。} 全负数组、单元素、和溢出、答案不能初始化为零。

\textbf{变体追问。} 若要求返回区间下标，同步记录重新开始位置和最优左右端。

\subsection{LeetCode 62 不同路径}

\textbf{题目模型。} 网格状态由上方和左方两种最后一步转移而来。

\textbf{推导重点。} 二维表可压缩为一维，因为当前行只依赖上一行同列和当前行左侧。

\textbf{边界。} 一行、一列、较大网格结果溢出、障碍物版本边界。

\textbf{变体追问。} 带障碍物时遇到障碍状态清零，第一行第一列要按阻断处理。

\subsection{LeetCode 63 不同路径 II}

\textbf{题目模型。} 障碍物让某些格子的路径数变为零。

\textbf{推导重点。} 滚动数组中遇到障碍直接把当前位置清零，否则累加左侧。

\textbf{边界。} 起点或终点是障碍、整行被截断、单行障碍。

\textbf{变体追问。} 若允许向上或向左走，图会有环，普通填表不再成立。

\subsection{LeetCode 64 最小路径和}

\textbf{题目模型。} 状态是到达当前格子的最小代价，最后一步来自上方或左方。

\textbf{推导重点。} 先初始化第一行和第一列，再处理内部，主转移保持简单。

\textbf{边界。} 单行、单列、权值为零、路径和溢出。

\textbf{变体追问。} 若允许四方向移动，变成最短路而不是普通网格 DP。

\subsection{LeetCode 70 爬楼梯}

\textbf{题目模型。} 最后一步来自前一阶或前两阶，是最小的重叠子问题模型。

\textbf{推导重点。} 滚动变量保存最近两个状态即可，不需要完整数组。

\textbf{边界。} n 为一、n 为二、题目是否允许零阶、结果范围。

\textbf{变体追问。} 若每次可走一到 k 阶，转移变成固定宽度窗口求和。

\subsection{LeetCode 72 编辑距离}

\textbf{题目模型。} 二维状态表示两个前缀之间的最少编辑操作。

\textbf{推导重点。} 末尾字符相同继承左上；不同则枚举插入、删除、替换。

\textbf{边界。} 空串、完全相同、完全不同、操作代价不同。

\textbf{变体追问。} 若只允许插入删除，问题退化到和最长公共子序列相关。

\subsection{LeetCode 121 买卖股票的最佳时机}

\textbf{题目模型。} 卖出日固定时，最佳买入日是此前最低价格。

\textbf{推导重点。} 扫描维护历史最低价和最大利润，一次交易无需复杂状态机。

\textbf{边界。} 单日价格、一直下跌、一直上涨、价格相同。

\textbf{变体追问。} 多次交易、冷冻期和手续费版本都要扩展成状态机。

\subsection{LeetCode 123 买卖股票 III}

\textbf{题目模型。} 最多两次交易，状态需要包含交易次数和持有状态。

\textbf{推导重点。} 维护 buy1、sell1、buy2、sell2 四个状态，按天更新。

\textbf{边界。} 价格为空、只够一次交易、两次交易间不能重叠。

\textbf{变体追问。} 最多 k 次交易是它的泛化，k 很大时退化为无限次交易。

\subsection{LeetCode 139 单词拆分}

\textbf{题目模型。} 状态表示前缀能否由字典单词拼出。

\textbf{推导重点。} 枚举最后一个单词的起点，左侧前缀可达且子串在字典中则当前可达。

\textbf{边界。} 空串、重复单词、长单词、substr 复制成本。

\textbf{变体追问。} 若要输出所有句子，需要 DFS 加记忆化或前驱列表。

\subsection{LeetCode 152 乘积最大子数组}

\textbf{题目模型。} 负数会让最大乘积和最小乘积互换。

\textbf{推导重点。} 同时维护以当前位置结尾的最大和最小乘积。

\textbf{边界。} 零、负数个数奇偶、单元素、乘积溢出。

\textbf{变体追问。} 如果要求返回区间，同步记录状态来源。

\subsection{LeetCode 198 打家劫舍}

\textbf{题目模型。} 每间房只有偷和不偷两种结局，偷当前就不能偷前一间。

\textbf{推导重点。} 滚动保存前一状态和前两状态。

\textbf{边界。} 空数组、单间、全零、金额很大。

\textbf{变体追问。} 环形房屋要拆成不含首和不含尾两个线性问题。

\subsection{LeetCode 221 最大正方形}

\textbf{题目模型。} 以当前位置为右下角的最大正方形由上、左、左上三个状态限制。

\textbf{推导重点。} 当前位置为一时取三者最小加一，否则为零。

\textbf{边界。} 全零、全一、单行单列、字符一和数字一不要混淆。

\textbf{变体追问。} 最大矩形则需要按行构造柱状图再用单调栈。

\subsection{LeetCode 300 最长递增子序列}

\textbf{题目模型。} 二次 DP 固定结尾，二分优化维护每个长度的最小结尾。

\textbf{推导重点。} tails 不是实际答案序列，而是未来扩展潜力表。

\textbf{边界。} 重复元素、严格递增和非递减区别、空数组。

\textbf{变体追问。} 若要还原路径，需要额外前驱数组，不能只看 tails。

\subsection{LeetCode 309 最佳买卖股票含冷冻期}

\textbf{题目模型。} 冷冻期让买入来源受昨天卖出状态限制。

\textbf{推导重点。} 状态机维护持有、刚卖出、休息三种日终状态。

\textbf{边界。} 空价格、一天、连续上涨、卖出后不能次日买入。

\textbf{变体追问。} 手续费版本可以在买入或卖出转移中扣费。

\subsection{LeetCode 714 买卖股票含手续费}

\textbf{题目模型。} 手续费改变交易收益，仍可用持有和不持有状态。

\textbf{推导重点。} 卖出时扣费或买入时扣费都可以，但不要重复扣。

\textbf{边界。} 手续费为零、价格震荡、小利润不足手续费。

\textbf{变体追问。} 这是状态机比局部贪心更稳的例子。

\subsection{LeetCode 1143 最长公共子序列}

\textbf{题目模型。} 状态表示两个前缀的最长公共子序列长度。

\textbf{推导重点。} 末尾相同取左上加一，不同取上方和左方最大。

\textbf{边界。} 空串、完全相同、无公共字符、子序列不是子串。

\textbf{变体追问。} 若要求还原序列，需要从表右下角反向追踪选择。

\subsection{背包模型}
下面只在这里讲一次共性：先写暴力基线，定位最贵的重复工作，再选择能维护不变量的数据结构。每道题卡只分析本题自己的建模、瓶颈、边界和变体，复习时不要把家族共性当成答案。

\begin{principlebox}{图示：背包模型推导链}
\begin{tabular}{@{}p{0.18\linewidth}p{0.74\linewidth}@{}}
暴力基线 & 枚举候选、完整模拟或逐个检查，先保证覆盖全部可能答案。\\
瓶颈定位 & 瓶颈是阶段相同、剩余容量相同的子问题反复出现。容量把未来选择压缩成一个数字，因此可以用表保存。\\
结构选择 & 0-1 背包每个物品只能用一次，一维压缩时容量倒序；完全背包物品可无限使用，容量正序。计数组合和计数排列的循环顺序不同。\\
不变量 & 处理到第 i 个物品后，\code{dp[cap]} 表示只使用前 i 类物品能达到的最优值或方案数。倒序是为了不在同一轮重复使用当前物品。\\
代码落地 & 容量通常与数值大小有关，所以复杂度是伪多项式。实现时先判断总和奇偶、目标范围和负数是否存在。\\
\end{tabular}
\end{principlebox}

\subsection{LeetCode 279 完全平方数}

\textbf{题目模型。} 完全平方数可以看成无限使用的物品，目标是凑出 n 的最少数量。

\textbf{推导重点。} 完全背包最少物品数，或把数字看成图节点做 BFS。

\textbf{边界。} n 为零或一、完全平方数本身、较大 n。

\textbf{变体追问。} 数学四平方定理可给更强解法，但面试 DP 更通用。

\subsection{LeetCode 322 零钱兑换}

\textbf{题目模型。} 金额是容量，硬币可无限使用，目标是最少硬币数。

\textbf{推导重点。} dp[value] 从 value 减 coin 的已知状态转移。

\textbf{边界。} 无法凑出、amount 为零、硬币面额大于目标、哨兵溢出。

\textbf{变体追问。} 零钱兑换 II 问组合数，循环语义不同。

\subsection{LeetCode 416 分割等和子集}

\textbf{题目模型。} 总和一半是背包容量，每个数只能使用一次。

\textbf{推导重点。} 0-1 背包可达性，一维容量必须倒序。

\textbf{边界。} 总和为奇数、目标为零、数值较大、复杂度依赖 target。

\textbf{变体追问。} 负数会破坏普通容量模型，需要重新建模。

\subsection{LeetCode 494 目标和}

\textbf{题目模型。} 正负号选择可代数转换为选若干数凑出特定正和。

\textbf{推导重点。} 若 S 加 target 为奇数或目标绝对值过大，直接无解。

\textbf{边界。} 零元素会让方案数翻倍、目标为负、容量为零。

\textbf{变体追问。} 也可用哈希 DP 保存所有可能和，但背包转换更高效。

\subsection{LeetCode 518 零钱兑换 II}

\textbf{题目模型。} 硬币无限使用，问组合数而不是最少数量。

\textbf{推导重点。} 外层遍历硬币、内层金额正序，保证组合按固定硬币顺序统计。

\textbf{边界。} amount 为零、无硬币、组合数较大、循环顺序反例。

\textbf{变体追问。} 若外层金额内层硬币，会统计排列数。

\subsection{贪心与交换证明}
下面只在这里讲一次共性：先写暴力基线，定位最贵的重复工作，再选择能维护不变量的数据结构。每道题卡只分析本题自己的建模、瓶颈、边界和变体，复习时不要把家族共性当成答案。

\begin{principlebox}{图示：贪心与交换证明推导链}
\begin{tabular}{@{}p{0.18\linewidth}p{0.74\linewidth}@{}}
暴力基线 & 枚举候选、完整模拟或逐个检查，先保证覆盖全部可能答案。\\
瓶颈定位 & 慢点是选择空间爆炸。若每一步都有很多候选，而某个排序规则或边界规则能安全丢掉大部分候选，就可能存在贪心。\\
结构选择 & 贪心需要找到可交换的局部最优：最早结束、最小右端、当前最远覆盖、当前最大收益或最小代价。排序只是手段，不是证明。\\
不变量 & 贪心不变量通常是当前方案在相同选择次数下不落后于任何其他方案，或者当前保留的资源最多、结束最早、右边界最小。\\
代码落地 & 实现时先把比较器写成明确的业务顺序。区间题注意起点相同时终点升序还是降序；覆盖题注意闭区间和半开区间的等号。\\
\end{tabular}
\end{principlebox}

\subsection{LeetCode 55 跳跃游戏}

\textbf{题目模型。} 扫描过程中只关心当前能到达的最远位置。

\textbf{推导重点。} 如果下标超过最远可达，任何之前路径都不能到达它；否则用它扩展边界。

\textbf{边界。} 第一个元素为零、数组长度一、恰好到达最后、远超最后。

\textbf{变体追问。} 求最少跳数时把可达区间看成 BFS 层。

\subsection{LeetCode 122 买卖股票 II}

\textbf{题目模型。} 允许多次交易且不能同时持有，所有正收益差都可以收集。

\textbf{推导重点。} 把长上涨段拆成每天买卖收益相同，因此累加正差值最优。

\textbf{边界。} 下降段、平台价格、只有一天、手续费不存在。

\textbf{变体追问。} 有手续费时不能简单累加正差，需要状态机或调整贪心阈值。

\subsection{LeetCode 169 多数元素}

\textbf{题目模型。} 超过一半的元素可以和其他元素两两抵消后仍然剩下。

\textbf{推导重点。} Boyer-Moore 投票维护候选和计数。

\textbf{边界。} 题目是否保证存在多数、计数归零、全相同。

\textbf{变体追问。} 若找超过三分之一的元素，需要维护两个候选。

\subsection{LeetCode 435 无重叠区间}

\textbf{题目模型。} 保留最多不重叠区间等价于删除最少区间。

\textbf{推导重点。} 按终点升序选择，结束越早给后续留下空间越多。

\textbf{边界。} 端点相接是否重叠、完全覆盖、相同终点、空列表。

\textbf{变体追问。} 用交换论证说明最早结束的选择不劣于其他第一选择。

\subsection{LeetCode 452 用最少数量的箭引爆气球}

\textbf{题目模型。} 一支箭的位置可以覆盖所有包含该位置的区间。

\textbf{推导重点。} 按右端点排序，当前箭放在最早结束气球的右端。

\textbf{边界。} 端点相同、负坐标、区间包含、闭区间边界。

\textbf{变体追问。} 它和无重叠区间是同一类最早结束贪心。

\subsection{LeetCode 763 划分字母区间}

\textbf{题目模型。} 每个片段必须覆盖其中所有字符的最后出现位置。

\textbf{推导重点。} 扫描维护当前片段最远右端，到达右端即可切分。

\textbf{边界。} 单字符、所有字符相同、字符多次跨段。

\textbf{变体追问。} 这是最早合法切分贪心，能最大化片段数量。

\subsection{区间、排序与合并}
下面只在这里讲一次共性：先写暴力基线，定位最贵的重复工作，再选择能维护不变量的数据结构。每道题卡只分析本题自己的建模、瓶颈、边界和变体，复习时不要把家族共性当成答案。

\begin{principlebox}{图示：区间、排序与合并推导链}
\begin{tabular}{@{}p{0.18\linewidth}p{0.74\linewidth}@{}}
暴力基线 & 枚举候选、完整模拟或逐个检查，先保证覆盖全部可能答案。\\
瓶颈定位 & 瓶颈是没有利用排序后相邻关系。区间按起点或终点排序后，很多远离的区间不可能再影响当前决策。\\
结构选择 & 合并按起点排序，调度按终点排序，覆盖题常用起点升序且终点降序。扫描时维护当前最远右端、结果尾区间或已选区间终点。\\
不变量 & 结果数组中的区间已经互不重叠且有序；当前扫描区间只可能影响结果最后一个区间或当前边界。\\
代码落地 & 比较器必须处理相同起点，否则覆盖和去重会错。区间端点可能溢出时用更大整数。闭区间和半开区间决定是否用小于等于。\\
\end{tabular}
\end{principlebox}

\subsection{LeetCode 56 合并区间}

\textbf{题目模型。} 排序后可能与当前区间相交的只会是结果尾部。

\textbf{推导重点。} 按起点排序，扫描时维护结果最后区间的右端点。

\textbf{边界。} 端点相接、完全覆盖、无重叠、空列表、输入未排序。

\textbf{变体追问。} 若区间是半开区间，端点相等不一定合并。

\subsection{LeetCode 57 插入区间}

\textbf{题目模型。} 新区间只会影响与它相交的一段连续区间。

\textbf{推导重点。} 先追加左侧完全不相交区间，再合并重叠段，最后追加右侧。

\textbf{边界。} 新区间在最前、最后、覆盖全部、刚好相邻、原列表为空。

\textbf{变体追问。} 若输入不是有序且不重叠，必须先排序或换模型。

\subsection{LeetCode 1288 删除被覆盖区间}

\textbf{题目模型。} 被覆盖要求另一区间起点不晚且终点不早。

\textbf{推导重点。} 按起点升序、终点降序排序，扫描维护最大右端。

\textbf{边界。} 同起点区间、完全覆盖、无覆盖、端点相等。

\textbf{变体追问。} 终点降序是关键，否则同起点小区间会先出现造成误判。

\subsection{回溯、枚举与剪枝}
下面只在这里讲一次共性：先写暴力基线，定位最贵的重复工作，再选择能维护不变量的数据结构。每道题卡只分析本题自己的建模、瓶颈、边界和变体，复习时不要把家族共性当成答案。

\begin{principlebox}{图示：回溯、枚举与剪枝推导链}
\begin{tabular}{@{}p{0.18\linewidth}p{0.74\linewidth}@{}}
暴力基线 & 枚举候选、完整模拟或逐个检查，先保证覆盖全部可能答案。\\
瓶颈定位 & 慢点在于大量分支在很早就已经不可能成功，仍然被继续展开。重复元素还会产生等价排列或组合。\\
结构选择 & 剪枝来自容量、剩余长度、排序后去重、合法性预检查和提前终止。组合题关注起点推进，排列题关注元素是否使用，分割题关注当前片段是否合法。\\
不变量 & 路径数组保存已经做出的选择；起点下标保证组合不回头；同层去重保证同一个位置层级不会选择相同值；已使用数组保证排列不重复使用同一元素。\\
代码落地 & 回溯代码虽然天然递归，但工程中要意识到深度。题目规模较小时递归可读性好；若要遵守严格工程栈限制，可以改成显式栈状态机。\\
\end{tabular}
\end{principlebox}

\subsection{LeetCode 78 子集}

\textbf{题目模型。} 每个元素都有选或不选两种选择，答案是所有路径节点。

\textbf{推导重点。} 回溯按起点向后枚举，避免不同顺序生成同一集合。

\textbf{边界。} 空数组、单元素、输出数量为二的 n 次方、顺序不要求。

\textbf{变体追问。} 也可用位掩码枚举所有子集，适合 n 较小的场景。

\subsection{LeetCode 79 单词搜索}

\textbf{题目模型。} 网格路径不能重复使用同一格，状态包含位置、匹配下标和访问标记。

\textbf{推导重点。} 剪枝包括首字符不匹配、剩余长度不足、字符频次不足和越界访问。

\textbf{边界。} 单字符单词、路径需要回退、重复字母、单元格不能复用。

\textbf{变体追问。} 若要查很多单词，应使用 Trie 合并搜索前缀。

\subsection{LeetCode 90 子集 II}

\textbf{题目模型。} 重复元素会让不同下标路径生成相同集合。

\textbf{推导重点。} 排序后在同一层跳过相同值，保留不同层使用重复值的可能。

\textbf{边界。} 全重复、空集、重复值相邻、去重条件写成同层而不是全局。

\textbf{变体追问。} 组合总和 II 使用同样的同层去重思想。

\subsection{LeetCode 131 分割回文串}

\textbf{题目模型。} 每次选择一个从当前位置开始的回文前缀。

\textbf{推导重点。} 预处理回文 DP 表，让回溯合法性检查为常数。

\textbf{边界。} 空串、单字符、全相同字符、重复分割路径。

\textbf{变体追问。} 若只求最少切割次数，可改成 DP 最短切分。

\subsection{LeetCode 216 组合总和 III}

\textbf{题目模型。} 从一到九中选 k 个数和为 n，每个数最多用一次。

\textbf{推导重点。} 按递增起点枚举，利用剩余个数和剩余和剪枝。

\textbf{边界。} n 太小、n 太大、k 为零、路径长度已经超过 k。

\textbf{变体追问。} 可预估剩余最小和最大和进一步剪枝。

\subsection{堆与动态极值}
下面只在这里讲一次共性：先写暴力基线，定位最贵的重复工作，再选择能维护不变量的数据结构。每道题卡只分析本题自己的建模、瓶颈、边界和变体，复习时不要把家族共性当成答案。

\begin{principlebox}{图示：堆与动态极值推导链}
\begin{tabular}{@{}p{0.18\linewidth}p{0.74\linewidth}@{}}
暴力基线 & 枚举候选、完整模拟或逐个检查，先保证覆盖全部可能答案。\\
瓶颈定位 & 慢点在于动态极值查询。我们只需要当前最极端元素，不需要整个集合完全有序。\\
结构选择 & 堆把插入和弹出极值控制在对数复杂度。Top K 用固定大小堆维护边界，合并多路序列用堆保存每一路当前头部，数据流中位数用双堆维护两半。\\
不变量 & 堆不变量不是全局有序，而是堆顶代表当前最优候选。若存在过期元素，需要在弹出时懒删除，保证真正使用堆顶前它仍然有效。\\
代码落地 & C++ 的 \code{std::priority_queue} 默认大顶堆。小顶堆用 \code{std::greater<T>}；结构体比较器要按谁优先级更低来写。堆中保存大对象时尽量保存下标或指针。\\
\end{tabular}
\end{principlebox}

\subsection{LeetCode 215 数组中的第 K 个最大元素}

\textbf{题目模型。} 只关心第 K 大，不需要完整排序。

\textbf{推导重点。} 大小为 K 的小顶堆保存当前前 K 大边界；快速选择可做到平均线性。

\textbf{边界。} K 为一、K 为 n、重复值、随机 pivot。

\textbf{变体追问。} 若数据流持续到来，堆方案比一次性快速选择更自然。

\subsection{LeetCode 295 数据流的中位数}

\textbf{题目模型。} 数据流需要动态维护较小一半和较大一半。

\textbf{推导重点。} 大顶堆保存较小一半，小顶堆保存较大一半，大小差不超过一。

\textbf{边界。} 偶数个元素、负数、重复值、平均值溢出。

\textbf{变体追问。} 若还要删除任意元素，需要延迟删除哈希表。

\subsection{LeetCode 347 前 K 个高频元素}

\textbf{题目模型。} 先统计频次，再只维护前 K 个候选。

\textbf{推导重点。} 小顶堆大小超过 K 就弹出最低频；桶排序利用频次上界。

\textbf{边界。} K 等于不同元素数、频次相同、结果顺序不要求。

\textbf{变体追问。} 若要按频次和字典序排序，比较器要同时处理两个键。

\subsection{LeetCode 621 任务调度器}

\textbf{题目模型。} 相同任务之间有冷却间隔，核心受最高频任务骨架限制。

\textbf{推导重点。} 可用大顶堆模拟，也可用最高频公式计算最短时间。

\textbf{边界。} 冷却为零、多个最高频、空闲被其他任务填满。

\textbf{变体追问。} 若任务有不同释放时间或执行时长，就变成更一般的调度问题。

\subsection{LeetCode 973 最接近原点的 K 个点}

\textbf{题目模型。} 距离只需比较平方，避免开方误差和成本。

\textbf{推导重点。} 大小为 K 的大顶堆保存当前最近 K 个点中最远者。

\textbf{边界。} K 为一、K 为全部、距离相同、坐标平方溢出。

\textbf{变体追问。} 快速选择可平均线性，但堆更适合数据流。

\subsection{链表与指针重连}
下面只在这里讲一次共性：先写暴力基线，定位最贵的重复工作，再选择能维护不变量的数据结构。每道题卡只分析本题自己的建模、瓶颈、边界和变体，复习时不要把家族共性当成答案。

\begin{principlebox}{图示：链表与指针重连推导链}
\begin{tabular}{@{}p{0.18\linewidth}p{0.74\linewidth}@{}}
暴力基线 & 枚举候选、完整模拟或逐个检查，先保证覆盖全部可能答案。\\
瓶颈定位 & 链表慢点在于不能随机访问。要找倒数位置、分组边界或交点，只能通过指针间距离、快慢指针或哈希记录访问过的节点。\\
结构选择 & 优化来自哨兵节点、快慢指针、前驱指针和局部反转。链表题的核心不是值交换，而是保持断开和接回的顺序正确。\\
不变量 & 必须明确每个指针指向哪一段：已经处理好的前缀头尾、当前待处理节点、下一段头节点。每次改 \code{next} 前先保存后继。\\
代码落地 & 节点通常由评测平台管理，代码不要随意 \code{delete} 输入节点。若创建新节点，要说明所有权。比较交点时比较地址，不比较值。\\
\end{tabular}
\end{principlebox}

\subsection{LeetCode 2 两数相加}

\textbf{题目模型。} 链表按低位到高位保存数字，本质是竖式加法而不是整数转换。

\textbf{推导重点。} 维护两个游标和进位，任一链未结束或进位非零都继续生成节点。

\textbf{边界。} 两链长度不同、最后进位、输入为零、不能用内置整数承载超长数字。

\textbf{变体追问。} 若链表高位在前，可以用栈反向处理或先反转链表。

\subsection{LeetCode 141 环形链表}

\textbf{题目模型。} 快慢指针在有环时必然相遇，无环时快指针先到空。

\textbf{推导重点。} 不用哈希表也能常数空间检测环。

\textbf{边界。} 空链、单节点自环、两节点环、尾部无环。

\textbf{变体追问。} 若要求环入口，用相遇后双指针同步走。

\subsection{LeetCode 142 环形链表 II}

\textbf{题目模型。} 相遇点到入口的距离关系来自快慢指针速度差。

\textbf{推导重点。} 相遇后一个指针回头，两个指针每次走一步，再次相遇就是入口。

\textbf{边界。} 入口在头节点、长尾短环、无环、单节点环。

\textbf{变体追问。} 也可用哈希集合，但空间更高。

\subsection{LeetCode 146 LRU 缓存}

\textbf{题目模型。} 哈希表负责按 key 定位，双向链表负责维护新旧顺序。

\textbf{推导重点。} 访问和更新都把节点移动到头部，容量满时删除尾部最旧节点。

\textbf{边界。} 容量为一、重复 put、get 不存在、更新已有值。

\textbf{变体追问。} C++ 可用 \code{std::list} 加迭代器，也可手写双向链表理解所有权。

\subsection{LeetCode 148 排序链表}

\textbf{题目模型。} 链表不能随机访问，适合归并排序而不是快速排序数组 partition。

\textbf{推导重点。} 自底向上迭代归并能避免递归栈，同时保持 O(n log n)。

\textbf{边界。} 空链、单节点、重复值、奇数长度、切断子链。

\textbf{变体追问。} 若把值复制到数组排序，会改变节点身份语义。

\subsection{LeetCode 160 相交链表}

\textbf{题目模型。} 相交是节点地址相同，不是节点值相同。

\textbf{推导重点。} 两个指针走完各自链表后切换到另一条链，走过总长度相同后相遇。

\textbf{边界。} 不相交、头节点相交、尾节点相交、长度差很大。

\textbf{变体追问。} 哈希集合更直观但空间更高。

\subsection{LeetCode 206 反转链表}

\textbf{题目模型。} 每次把当前节点的 next 指向已经反转好的前缀头。

\textbf{推导重点。} 先保存后继，再改指针，再推进 previous 和 current。

\textbf{边界。} 空链、单节点、两个节点、不要丢失剩余链。

\textbf{变体追问。} K 组反转是在这段局部反转外增加分组边界。
